The CosMx application and locked references are Table~\ref{tab:materials}. The tumor application uses CosMx basal-cell carcinoma cells from four patients~\cite{Andreatta2024zenodo,Yerly2022,Yerly2025preprint}. The export comprises cell barcodes, gene features, per-cell metadata, and the counts matrix. Metadata carry Patient\_ID, Condition, and slide name. Patients A and B are labelled Ulcerated nodular. Patients C and D follow a wound time course: Baseline, Unwound, and Wound. The four-patient export contains 232{,}802 cells. Restriction to the eight shared types leaves 216{,}949 typed cells: Patient A, 20{,}730; Patient B, 41{,}728; Patient C, 90{,}487; Patient D, 64{,}004.

The tumor application uses CosMx BCC, Cancer.cells versus Normal.Kerat~\cite{He2022,Andreatta2024zenodo}. An independent CosMx NSCLC pair uses pooled tumor versus fibroblast (cosine \(0.648151\)) on the NanoString NSCLC FFPE release, whose 960-plex prototype panel is described in the CosMx platform report~\cite{He2022,Sun2025spatial}. A CosMx colorectal section uses pooled epithelium versus fibroblast (cosine \(0.505826\))~\cite{Crowell2025preprint}. A CosMx HCC sample uses pooled tumor versus stellate cells (cosine \(0.648597\)) on the NanoString human liver FFPE release~\cite{Sun2025spatial}. A CosMx PDAC library uses tumor versus CAF (cosine \(0.861599\))~\cite{Pei2025}. The same frozen cosine is then read on three additional libraries: HNSCC openST (Tumor versus Tumor\_Keratin\_Pearl)~\cite{Schott2024}, breast Xenium (Invasive\_Tumor versus Prolif\_Invasive\_Tumor)~\cite{Janesick2023}, and an orthogonal Xenium FLEX reference on the same 2864 Xenium spots (cosine \(0.987505\))~\cite{Janesick2023}. Full-\(n\) confirmation at \(t=0\) uses 6971 openST spots, 2864 Xenium spots, and 5686 CosMx spots. Type-pair identities and \(c^\star=0.80\) are frozen before any sweep.

The C-from-D composition (overall RMSE \(0.1149\), type-level PCC \(0.7708\), spot-level PCC \(0.8773\)) is recorded in Table~\ref{tab:materials}.
